\section*{Chapter \ref{ch:b3:atmospheric-dynamics} --- Atmospheric Dynamics}

\begin{solution}{exo:b3:atmospheric-dynamics:1}
(a)~$\zeta = 2\omega = 2\times2\pi/86\,400 = \qty{1.45e-4}{s^{-1}}$.
(b)~$f = \num{7.29e-5}$ and \qty{1.26e-4}{s^{-1}}: the vortex has about
twice the \omterm{def:b3:atmospheric-dynamics:vorticity-kinds}{planetary vorticity} of \ang{30} N and a little more than that of
\ang{60} N.
\end{solution}

\begin{solution}{exo:b3:atmospheric-dynamics:2}
$\zeta = f_{30} - f_{40} = \num{7.29e-5} - \num{9.37e-5} =
\qty{-2.08e-5}{s^{-1}}$: anticyclonic.
\end{solution}

\begin{solution}{exo:b3:atmospheric-dynamics:3}
$k = 2\pi/(5\times10^6) = \qty{1.257e-6}{m^{-1}}$; $c = -\beta/k^2 =
-\num{1.62e-11}/\num{1.58e-12} = \qty{-10.3}{m/s}$: westward.
\end{solution}

\begin{solution}{exo:b3:atmospheric-dynamics:4}
$\beta = 2\Omega\cos\ang{60}/a = \qty{1.14e-11}{m^{-1}.s^{-1}}$;
$\lambda_s = 2\pi\sqrt{10/\num{1.14e-11}} = \qty{5870}{km}$; the circle,
\qty{20015}{km}, holds $n_s = 3.4$ waves.
\end{solution}

\begin{solution}{exo:b3:atmospheric-dynamics:5}
$q = f/h$ is conserved: over the range $\zeta = f(8/10) - f =
\qty{-2.06e-5}{s^{-1}}$, anticyclonic; over the basin $f(12/10) - f =
\qty{+2.06e-5}{s^{-1}}$, cyclonic. Where columns are stretched, in the lee
of the range and over the basin, cyclonic vorticity is created: there
depressions tend to form.
\end{solution}

\begin{solution}{exo:b3:atmospheric-dynamics:6}
$k = 5/(a\cos\ang{45}) = \qty{1.11e-6}{m^{-1}}$, $\beta/k^2 =
\qty{13.1}{m/s}$: $c = \qty{1.9}{m/s}$, $c_g = \qty{28.1}{m/s}$. The energy
covers \qty{5000}{km} in \num{2.1}~days, the pattern in \num{31}~days.
\end{solution}

\begin{solution}{exo:b3:atmospheric-dynamics:7}
(a)~$k^2 = \beta/U - l^2 = \num{1.079e-12} - \num{2.74e-13} =
\qty{8.05e-13}{m^{-2}}$, $k = \qty{8.97e-7}{m^{-1}}$: $\lambda =
\qty{7000}{km}$ and $n = 4.0$, fewer waves than with $l = 0$. (b)~$l^2 =
\num{1.097e-12} > \beta/U = \num{1.079e-12}$: $k^2$ would be negative.
\end{solution}

\begin{solution}{exo:b3:atmospheric-dynamics:8}
$L_d = 0.012\times10^4/\num{1.031e-4} = \qty{1164}{km}$; $\sigma =
0.31\times\num{1.031e-4}\times\num{2e-3}/0.012 = \qty{5.3e-6}{s^{-1}}$, an
$e$-folding time of \num{2.2}~days; wavelength $2\pi\times1164/1.61 =
\qty{4540}{km}$.
\end{solution}

\begin{solution}{exo:b3:atmospheric-dynamics:9}
$[v^*T^*] = (VT_0/2)\cos\delta$: \qty{25}{K.m/s}, \qty{12.5}{K.m/s} and 0;
heat fluxes \num{25100}, \num{12600} and \qty{0}{W/m^2}. When the
temperature is a quarter wavelength out of phase with the wind, as much
warm air goes south as north: no transport.
\end{solution}

\begin{solution}{exo:b3:atmospheric-dynamics:10}
Linearising about $U$ with $\zeta' = \nabla^2\psi'$ and $v' =
\partial_x\psi'$: $(\partial_t + U\partial_x)\nabla^2\psi' + \beta\partial_x
\psi' = 0$. With $\psi' \propto \mathrm e^{\mathrm i(kx + ly - \omega t)}$,
$(\omega - Uk)K^2 = -\beta k$, so $\omega = Uk - \beta k/K^2$. Relative to
the air the phase speed is $c - U = -\beta/K^2 < 0$ for every wave.
\end{solution}

\begin{solution}{exo:b3:atmospheric-dynamics:11}
$\varphi_H = [5\times9.81\times12\,000\times0.25/(3\times(\num{7.292e-5})^2
\times(\num{6.371e6})^2)]^{1/2} = \qty{0.477}{rad}$, \ang{27.3}.
$u_M = \Omega a\varphi_H^2 = \qty{106}{m/s}$; exactly, $\Omega a\sin^2
\varphi_H/\cos\varphi_H = \qty{110}{m/s}$: two and a half times the
observed jet.
\end{solution}

\begin{solution}{exo:b3:atmospheric-dynamics:12}
(a)~Across the band the interface is displaced by $\eta = sy$, $-L/2 \le y
\le L/2$, from its flat position. A layer of reduced density contrast
displaced by $\eta$ stores $\tfrac12\rho g'\eta^2$ per unit area; averaged
over the band, $\langle\eta^2\rangle = s^2L^2/12$. (b)~$g' = 0.327$:
\num{1.5e5} and \qty{6.1e5}{J/m^2}, against $\tfrac12\times0.7\times15^2
\times10^4 = \qty{7.9e5}{J/m^2}$ of kinetic energy: comparable.
\end{solution}

\begin{solution}{pb:b3:atmospheric-dynamics:1}
\textbf{1.} $f_{45} = \qty{1.031e-4}{s^{-1}}$, $f_{50} =
\qty{1.117e-4}{s^{-1}}$; $\zeta = f_{45} - f_{50} = \qty{-8.6e-6}{s^{-1}}$.

\textbf{2.} $R = V/|\zeta| = 15/\num{8.6e-6} = \qty{1750}{km}$; clockwise,
back towards the south.

\textbf{3.} $\zeta = f_{45}(9.5/10) - f_{50} = \qty{-1.38e-5}{s^{-1}}$, more
anticyclonic.

\textbf{4.} $\zeta = f_{45} - f_{40} = \qty{+9.4e-6}{s^{-1}}$, cyclonic: it
turns back north. Displaced either way, the parcel acquires the vorticity
that turns it back: a restoring effect, and an oscillation.

\textbf{5.} $\beta = \qty{1.62e-11}{m^{-1}.s^{-1}}$; the circle is
\qty{28306}{km} long.

\textbf{6.} $n = 3$ to 8: wavelengths 9435, 7076, 5661, 4718, 4044 and
\qty{3538}{km}; $c = -21.5$, $-5.5$, 1.9, 5.9, 8.3 and \qty{9.9}{m/s}; $c_g
= 51.5$, 35.5, 28.1, 24.1, 21.7 and \qty{20.1}{m/s}.

\textbf{7.} $n = 3$ and 4 move west, $n \ge 5$ east. The wave $n = 5$ takes
$10^6/1.86 = \qty{5.4e5}{s}$, \num{6.2}~days, to move \qty{1000}{km}.

\textbf{8.} Downstream, east: at $c_g = \qty{24.1}{m/s}$ the energy covers
\qty{5000}{km} in \num{2.4}~days.

\textbf{9.} $\lambda_s = 2\pi\sqrt{15/\num{1.62e-11}} = \qty{6050}{km}$,
$n_s = 4.7$.

\textbf{10.} $n_s = 5.7$ for \qty{10}{m/s} and 4.1 for \qty{20}{m/s}: a
stronger westerly favours longer stationary waves, fewer round the
hemisphere.

\textbf{11.} $n_s = 5.8$ at \ang{35} N and 3.4 at \ang{55} N.

\textbf{12.} $\zeta = f(8/10) - f = \qty{-2.06e-5}{s^{-1}}$: squeezing over
the mountains makes the column anticyclonic, stretching in the lee
cyclonic. The mountains impose a fixed pattern of vorticity on the flow,
a forcing that does not move.

\textbf{13.} A fixed forcing produces a steady response; the response is
strongest when the forced pattern is itself a free wave that would stand
still, the resonance at $\lambda_s$.

\textbf{14.} $\Lambda = 25/10^4 = \qty{2.5e-3}{s^{-1}}$;
$|\partial T/\partial y| = \Lambda fT/g = \num{2.5e-3}\times\num{1.031e-4}
\times250/9.81 = \qty{6.6}{K}$ per \qty{1000}{km}.

\textbf{15.} $\mathrm{Ri} = 0.011^2/(\num{2.5e-3})^2 = 19$: far above 0.25,
no shear turbulence.

\textbf{16.} $L_d = 0.011\times10^4/\num{1.031e-4} = \qty{1067}{km}$.

\textbf{17.} $2\pi\times1067/1.61 = \qty{4160}{km}$; $28306/4160 = 6.8$
waves.

\textbf{18.} $2\pi\times1067/2.40 = \qty{2790}{km}$.

\textbf{19.} $\sigma = 0.31\times\num{1.031e-4}\times\num{2.5e-3}/0.011 =
\qty{7.3e-6}{s^{-1}}$, \num{0.63}~day$^{-1}$.

\textbf{20.} $\exp(3\times0.627) = 6.6$.

\textbf{21.} The storm, \qty{4160}{km}, is shorter than the stationary
wavelength, \qty{6050}{km}: as a Rossby wave it is carried east by the
westerly, and the storms travel east through the slower or standing long
waves.

\textbf{22.} Westward with height: only then is the poleward-moving air
warmer than the equatorward-moving air, so that heat goes poleward and the
\omterm{def:b3:atmospheric-dynamics:ape}{available potential energy} is released.

\textbf{23.} The fastest-growing disturbance grows by a factor $e$ every
\textbf{\num[text-series-to-math=true]{1.6}~days} (\num{1.59}~days by the
calculation).
\end{solution}
